%% ma: L10-B02-0011
%% lop: 10
%% chuong: 1
%% bai: 2
%% dang: 10.2.7
%% loai: TN4
%% muc: VD
%% dap_an: "B"
%% nguon: "Ảnh giáo viên gửi 10/10/2026 (ThS. Nguyễn Hoàng Việt, Chương 2 Tập hợp), A-Đề 1, Phần 1 Câu 9"
%% kien_thuc: "Bài 2 (tập hợp và các phép toán trên tập hợp)"
%% trang_thai: chinh-thuc
%% ghi_chu: "Đề gốc không có đáp án; tự giải được m=9, chỉ có phương án B chứa 9"
\begin{ex}
Cho tập hợp $A=(0;+\infty)$ và $B=\{x\in\mathbb R\mid mx^2-6x+m-8=0\}$, $m$ là tham số. Biết rằng tập $B$ có đúng hai tập con và $B\subset A$. Kết luận nào sau đây là đúng về giá trị của tham số $m$?
\choice{$m\in(0;9)$}{\True $m\in(0;10)$}{$m\in(0;8)$}{$m\in(0;7)$}
\loigiai{$B$ có đúng hai tập con nên $B$ có đúng một phần tử, tức phương trình $mx^2-6x+m-8=0$ có đúng một nghiệm thực.
Nếu $m=0$: $-6x-8=0\Leftrightarrow x=-\dfrac43\notin A$, loại.
Nếu $m\ne0$: $\Delta'=9-m(m-8)=0\Leftrightarrow m^2-8m-9=0\Leftrightarrow m=9$ hoặc $m=-1$. Nghiệm kép $x=\dfrac3m$: $m=9$ cho $x=\dfrac13>0$ (nhận); $m=-1$ cho $x=-3\notin A$ (loại).
Vậy $m=9$, nên $m\in(0;10)$.}
\end{ex}
